
\SGARefTarget{sga3:ix:section:8}
\section*{8. Subgroups, Quotient Groups, and Extensions of Groups of
Multiplicative Type over a Field\textsuperscript{*}}
\addcontentsline{toc}{section}{8. Subgroups, Quotient Groups, and Extensions
of Groups of Multiplicative Type over a Field}
\begingroup
\renewcommand{\thefootnote}{*}
\footnotetext{Added in July 1969. This afterthought section is independent
of Sections \SGARefLink{sga3:ix:section:3}{3} through \SGARefLink{sga3:ix:section:7}{7}.}
\endgroup

\medskip
\noindent\SGARefTarget{sga3:ix:scholium:8:0}\textbf{Scholium 8.0.}
\footnote{Editor's note: this scholium has been added.} ---
\textit{Let \(k\) be a field and let \(H\) be an affine \(k\)-group
scheme. Write \(k[H]\) for its affine algebra and
\[
  X(H)=\operatorname{Hom}_{k\text{-}\mathrm{gr}}
    (H,\mathbf G_{m,k})
\]
for the group of characters of \(H\). By the independence lemma for
characters, \(X(H)\) is a linearly independent subset of \(k[H]\).
Consequently, \(H\) is diagonalizable if and only if \(X(H)\) spans
\(k[H]\) as a \(k\)-vector space.}

\medskip
\noindent\SGARefTarget{sga3:ix:proposition:8:1}\textbf{Proposition 8.1.} ---
\textit{Let \(G\) be a diagonalizable group scheme (respectively, a group
scheme of multiplicative type) over a field \(k\). Then, for
every subgroup scheme \(H\) of \(G\), the schemes \(H\) and \(G/H\) are
diagonalizable (respectively, of multiplicative type).}

Definition \SGARefLink{sga3:ix:definition:1:1}{1.1} immediately reduces
the assertion to the case without the parenthetical qualification. An
easy passage to the limit, using VI\textsubscript{B} \SGARefLink{prop:sga3-VIB-11.13}{11.13} and \SGARefLink{sga3:viii:theorem:3:1}{VIII 3.1},
reduces us to the case in which \(G\) is of finite type over
\(k\).
By VI\textsubscript{B} \SGARefLink{thm:sga3-VIB-11.16}{11.16}, one can find a finite
family of nonzero elements
\SGARefTarget{sga3:ix:eq:8:1:1}
\begin{equation}
  \tag{8.1.1}
  f_i=\sum_m a_{im}m,\qquad a_{im}\in k,
\end{equation}
of the affine ring \(k(M)\) of \(G=D_k(M)\), such that the points of
\(H\), with values in an arbitrary \(k\)-algebra \(k'\), are precisely
the points \(g\in G(k')\) for which
\SGARefTarget{sga3:ix:eq:8:1:2}
\begin{equation}
  \tag{8.1.2}
  \tau_g f_i=\lambda_i(g)f_i,
  \qquad \lambda_i(g)\in k',
\end{equation}
where \(\tau_g\) denotes translation by \(g\). But
\SGARefTarget{sga3:ix:eq:8:1:3}
\begin{equation}
  \tag{8.1.3}
  \tau_g f_i=\sum_m a_{im}\chi_m(g)m,
\end{equation}
where \(\chi_m:G\to\mathbf G_{m,k}\) is the character associated with
\(m\). Thus \SGARefLink{sga3:ix:eq:8:1:2}{(8.1.2)} is equivalent to
\SGARefTarget{sga3:ix:eq:8:1:4}
\begin{equation}
  \tag{8.1.4}
  \chi_m(g)=\chi_{m'}(g)
  \qquad\text{if }m,m'\in Z_i,
\end{equation}
where \(Z_i\) denotes the set of \(m\in M\) for which \(a_{im}\neq0\).
\SGARefLink{sga3:ix:eq:8:1:4}{The same relation} may be written
\[
  \chi_{m'-m}(g)=1
  \qquad\text{if }m,m'\in Z_i.
\]

Let \(N\) be the subgroup of \(M\) generated by all the elements
\(m'-m\), where \(i\) varies and \(m,m'\in Z_i\). It follows from the
definition of \(D_k(M)\simeq G\) that \(H\) identifies with
\(D_k(M/N)\). Hence \SGARefLink{sga3:viii:theorem:3:1}{VIII 3.1} identifies \(G/H\) with \(D_k(N)\).
\hfill\(\square\)

\medskip
\noindent\SGARefTarget{sga3:ix:proposition:8:2}\textbf{Proposition 8.2.}
\footnote{Editor's note: both the statement and its proof have been
expanded.} ---
\textit{Let \(k\) be a field; let \(H\) and \(K\) be \(k\)-group schemes
of multiplicative type and finite type; and let \(G\) be a \(k\)-group
scheme fitting into an exact sequence
\SGARefTarget{sga3:ix:eq:8:2:0}
\[
  1\longrightarrow H\longrightarrow G\longrightarrow K\longrightarrow 1
\]
(which implies that \(G\) is of finite type over \(k\)).}

\SGARefTarget{sga3:ix:proposition:8:2:item:i}\textit{(i) If \(G\) is commutative or \(K\) is connected, then \(G\) is
of multiplicative type.}

\SGARefTarget{sga3:ix:proposition:8:2:item:ii}\textit{(ii) If \(K\) and \(H\) are diagonalizable and \(K\) is a torus,
then \(G\) is diagonalizable.}

For the proof of \SGARefLink{sga3:ix:proposition:8:2:item:i}{(i)}, the reader is referred to \SGARefLink{sga3:xvii:proposition:7:1:1}{XVII 7.1.1}, of which
\SGARefLink{sga3:ix:proposition:8:2:item:i}{8.2(i)} is a special case. The case
of a field is treated in part \SGARefLink{sga3:xvii:proposition:7:1:1:item:i}{(i)} of the proof of \SGARefLink{sga3:xvii:proposition:7:1:1}{XVII 7.1.1}, without
using the results of the later exposés.

For \SGARefLink{sga3:ix:proposition:8:2:item:ii}{(ii)}, now suppose that \(K\) and \(H\) are diagonalizable:
\[
  K\simeq D_k(M)
  \qquad\text{and}\qquad
  H\simeq D_k(N),
\]
and suppose that \(G\) is of multiplicative type, as follows from \SGARefLink{sga3:ix:proposition:8:2:item:i}{(i)}
when \(K\) is a torus. By \SGARefLink{sga3:x:proposition:1:4}{X 1.4}, \(G\) is isotrivial over \(k\); that is,
there is a finite separable extension \(k'/k\) such that
\[
  G'=G\times_k k'
\]
is diagonalizable. Hence \(G'=D_{k'}(E)\) for a commutative group \(E\),
and there is an exact sequence
\SGARefTarget{sga3:ix:eq:8:2:1}
\begin{equation}
  \tag{1}
  0\longrightarrow D_{k'}(N)\longrightarrow D_{k'}(E)
    \longrightarrow D_{k'}(M)\longrightarrow 0.
\end{equation}
By \SGARefLink{sga3:viii:theorem:3:1}{VIII 3.1} and \SGARefLink{sga3:ix:theorem:3:2}{3.2}, \(M\) is therefore a subgroup of \(E\), and one has
an exact sequence
\SGARefTarget{sga3:ix:eq:8:2:2}
\begin{equation}
  \tag{2}
  0\longrightarrow M\longrightarrow E\longrightarrow N\longrightarrow 0.
\end{equation}

For a given extension \(E_0\) of \(N\) by \(M\), consider the
diagonalizable \(k\)-group \(G_0=D_k(E_0)\). The \(k\)-group functor
\(A\) of automorphisms of the extension
\SGARefTarget{sga3:ix:eq:8:2:3}
\begin{equation}
  \tag{3}
  1\longrightarrow H\longrightarrow G_0\longrightarrow K\longrightarrow 1,
\end{equation}
that is, the group subfunctor of
\(\operatorname{Aut}_{k\text{-}\mathrm{gr}}(G_0)\) whose points on a
\(k\)-prescheme \(T\) are the
\[
  \phi\in\operatorname{Aut}_{T\text{-}\mathrm{gr}}(G_T)
\]
inducing the identity on \(H_T\) and \(K_T\), identifies with
\(\operatorname{Hom}_{k\text{-}\mathrm{gr}}(K,H)\). By \SGARefLink{sga3:viii:corollary:1:5}{VIII 1.5}, this is
the constant \(k\)-group with value
\[
  L=\operatorname{Hom}_{\mathrm{gr}}(N,M).
\]

It follows that the extensions \(G\) of \(K\) by \(H\) which become
isomorphic to \SGARefLink{sga3:ix:eq:8:2:3}{(3)} over a separable closure
\(k_s\) of \(k\) are classified in the same way as the \(k\)-torsors,
for the étale topology, under the constant group \(L_k\). Put
\(\Gamma=\operatorname{Gal}(k_s/k)\). Since \(L\) is a trivial
\(\Gamma\)-module, these torsors are classified by
\SGARefTarget{sga3:ix:eq:8:2:4}
\[
  H^1_{\mathrm{\acute et}}(k,L_k)
  =H^1(\Gamma,L)
  =\operatorname{Hom}_{\mathrm{gr.\,top.}}(\Gamma,L).
\]

If \(K\) is moreover a torus, then \(M\) is torsion-free, and hence so
is \(L\). Therefore
\[
  \operatorname{Hom}_{\mathrm{gr.\,top.}}(\Gamma,L)=0,
\]
and every extension of \(K\) by \(H\) is already diagonalizable.

\medskip
\noindent\SGARefTarget{sga3:ix:remark:8:3}\textbf{Remark 8.3.} ---
\textit{As already noted in the proof of
\SGARefLink{sga3:ix:proposition:8:2}{8.2}, \SGARefLink{sga3:ix:proposition:8:2:item:i}{the first assertion} is stated
and proved over an arbitrary base in \SGARefLink{sga3:xvii:proposition:7:1:1}{XVII 7.1.1}. \SGARefLink{sga3:ix:proposition:8:2:item:ii}{The second assertion}
generalizes, with essentially the same proof, to the case of an integral
normal base scheme---or, more generally, a geometrically unibranch
one---by using \SGARefLink{sga3:x:lemma:5-13}{X 5.13} below.}

\textit{Finally, \SGARefLink{sga3:ix:proposition:8:1}{8.1} may also be
regarded as a corollary of the distinctly less trivial result \SGARefLink{sga3:ix:theorem:6:8}{6.8}, which
is likewise valid over an arbitrary base scheme.}

\SGARefTarget{sga3:ix:bibliography}
\section*{Bibliography}
\addcontentsline{toc}{section}{Bibliography}
\footnote{Editor's note: additional references cited in this exposé.}

\noindent \SGARefTarget{sga3:ix:bib:bens}[BEns] N. Bourbaki, \textit{Théorie des ensembles},
Chapters I--IV, Hermann, 1970.
